\chapter{The Problem}
\label{ch:problem}

\section{Dead inventory in the thinking status line}

Every time an AI coding agent pauses to ``think,'' it renders a status line ---
a spinner, a shimmer, a few words of ephemeral text --- that conveys nothing of
commercial value and earns nothing for anyone. With 85--92\% of developers now
using AI coding tools and the most engaged among them spending 4--10 hours a
day inside them, this adds up to an enormous, continuously refreshed pool of
captive, full-attention wait-state moments \citep{masterplan_2026}. The
moment is unusually valuable as an attention surface precisely because it is
\emph{not} a sidebar a user scrolls past: the developer is waiting, watching,
and briefly idle. Today that attention is simply discarded. It is dead
inventory.

The scale of the underlying surface is not speculative. Claude Code ---
Devorix's first target tool --- grew from \$0 to a \$2.5B run-rate in nine
months, and Cursor reached $\sim$\$2B annualized revenue with over a million
paying customers \citep{masterplan_2026}. \emph{(Verified fact.)} Real,
large, fast-growing engagement is flowing through the exact surface where this
wait-state moment occurs.

\section{Developers earn nothing from the time they spend}

Developers spend that wait time --- and the broader hours they pour into AI
coding tools --- without any financial participation in the value their
attention represents. For the large under-monetized developer base this is a
structural mismatch: GitHub Sponsors supported 49{,}148 developers as of March
2026, yet 60\% of surveyed open-source maintainers describe themselves as
unpaid hobbyists and only 13\% as professional or paid
\citep{customerproduct_2026}. \emph{(Verified fact.)} A large population is
structurally primed to want incremental income attached to tools they already
use.

We are deliberately honest about the shape of this problem, following the
plan's own VC-lens validation \citep{masterplan_2026}. This is
\textbf{a monetization opportunity layered onto an existing attention pool, not
a pain point developers are actively trying to solve}. Developers are not
asking for this. The closest behavioural precedent --- cashback browser
extensions that retain users because the reward appears ambiently inside an
existing workflow rather than requiring a separate action --- is encouraging
for adoption but is consumer shopping, not developer tooling
\citep{customerproduct_2026}. \emph{(Industry estimate.)} The pain that is
genuinely severe sits on the \textbf{advertiser} side: developer-tool marketers
have shockingly few high-trust placement options built for the agentic-coding
moment, with Carbon Ads (audited at \$1.60 CPM, selling to Okta, GitLab and
MongoDB) essentially the only mature one, and it was never built for this
moment specifically \citep{masterplan_2026}.

\section{A working competitor exists --- but structurally cannot pay India}

The wait-state advertising mechanic is not hypothetical. Kickbacks.ai launched
on 2026-06-11 and its launch tweet passed 5.5~million views in 24~hours; a
direct clone, IdleAds.dev, appeared within 24~hours of it
\citep{masterplan_2026}. \emph{(Verified fact.)} The category has real
top-of-funnel virality. But it has a structural hole that defines Devorix's
wedge:

\textbf{Kickbacks --- and the Stripe-based model generally --- structurally
cannot pay out to India.} As of June~2026, Kickbacks had not yet activated real
payouts; its Stripe Connect integration was reported ``nearing completion''
\citep{masterplan_2026}. The reason is independently re-verified beyond a
single GitHub issue, against Stripe's own support documentation: Stripe has
been invite-only for new India accounts since May~2024 due to RBI regulatory
changes, and Connect Express onboarding for India-based recipients is
restricted \citep{stripe_india_2026}. Confidence on this wedge is
\textbf{High}. The same Stripe-India gap that is slowing the category leader is
exactly what Devorix is built around: native, automated, GST-ready UPI payouts
via RazorpayX, which supports automated UPI/IMPS/NEFT/RTGS with 24/7 instant
settlement \citep{razorpayx_2026}.

We are careful not to overclaim here. IdleAds.dev already does UPI
\emph{manually}, so the payout-rail edge is \emph{contested, not exclusive} ---
the real differentiator is native, automated, GST-ready UPI done properly,
not merely ``we pay India'' \citep{masterplan_2026}.

\section{Why this is a real, timely problem}

The timing is right, narrowly, on three converging facts
\citep{masterplan_2026}:

\begin{enumerate}
  \item \textbf{Adoption just crossed the majority case.} AI coding tools moved
        from early-adopter to mainstream (85--92\% of developers), so the
        attention pool the mechanic monetizes is now genuinely large.
        \emph{(Verified fact.)}
  \item \textbf{India is the fastest-growing developer base anywhere} --- 27M
        on GitHub in 2026, projected 57.5M by 2030, with India among the
        highest-adoption regions for AI coding tools --- yet it is the one
        large market the dominant Stripe-based payout model cannot serve.
        \emph{(Verified fact / India-specific adoption: Industry estimate.)}
  \item \textbf{The bottleneck is wide open.} Zero verified paying advertisers
        exist anywhere in this category; the real constraint is demand-side
        sales, which every competitor is failing at while racing the same
        shallow mechanic \citep{masterplan_2026}. \emph{(High confidence ---
        absence of evidence corroborated across a 3-day pulse-check.)}
\end{enumerate}

The honest read from the plan's VC-lens validation is preserved here without
softening: the window is short and the moat is shallow --- IdleAds cloned
Kickbacks' mechanic within 24~hours --- so timing favours moving fast on the
demand side, where the real, severe, currently-unmet need lives, rather than on
product polish on the supply side, where developers are mildly interested at
best \citep{masterplan_2026}. The problem Devorix addresses is therefore
real and timely, but its sharpest edge is the advertiser's missing high-trust
India-payable placement, not a developer crying out to be paid.
